% Рисунок с доски: отклик модели Кельвина — Фойгта на импульс нагрузки
% S = S_0 H(t) - S_0 H(t - t*).
\begin{figure}[htbp]
\centering
\begin{tikzpicture}[>={Latex[length=4pt]}, xscale=1.0]
  % нагрузка
  \begin{scope}
    \draw[->] (0,0) -- (6.3,0) node[anchor=north east] {$t$};
    \draw[->] (0,0) -- (0,2.2) node[anchor=south east] {$\Sdev$};
    \draw[very thick, red!60!black] (0,0) -- (0,1.6) -- (2.5,1.6) -- (2.5,0) -- (6.0,0);
    \node[anchor=north, font=\footnotesize] at (2.5,-0.05) {$t^*$};
    \node[anchor=east, font=\footnotesize] at (-0.05,1.6) {$\Sdev_0$};
  \end{scope}
  % отклик
  \begin{scope}[xshift=7.6cm]
    \draw[->] (0,0) -- (6.3,0) node[anchor=north east] {$t$};
    \draw[->] (0,0) -- (0,2.2) node[anchor=south east] {$\Lambda$};
    \draw[very thick, blue!60!black, domain=0:2.5, samples=80, smooth]
      plot (\x, {1.6*(1-exp(-\x/1.0))});
    \draw[very thick, blue!60!black, domain=2.5:6, samples=80, smooth]
      plot (\x, {1.6*(exp(-(\x-2.5)/1.0)-exp(-\x/1.0))});
    \draw[dashed, gray] (2.5,0) -- (2.5,{1.6*(1-exp(-2.5))});
    \node[anchor=north, font=\footnotesize] at (2.5,-0.05) {$t^*$};
  \end{scope}
\end{tikzpicture}
\caption{Импульс нагрузки и отклик по формуле Стилтьеса:
$\dev(t) = \Sdev_0\bigl(\Pi(t) - \Pi(t-t^*)\bigr) = \Sdev_0\Lambda(t)$.
После снятия нагрузки деформация убывает по экспоненте, а не исчезает мгновенно.}
\label{fig:creep-pulse}
\end{figure}
